\documentclass[11pt,a4paper]{article}
\usepackage[margin=23mm]{geometry}
\usepackage[T1]{fontenc}
\usepackage[utf8]{inputenc}
\usepackage{lmodern,microtype,amsmath,amssymb,booktabs,tabularx,array}
\usepackage{graphicx,xcolor,listings,enumitem}
\usepackage[section]{placeins}
\usepackage[colorlinks=true,linkcolor=blue!50!black,urlcolor=blue!50!black]{hyperref}
\usepackage{fancyhdr}
\pagestyle{fancy}
\fancyhf{}
\lhead{Statistical learning: worked notebook explanations}
\rfoot{\thepage}
\setlength{\headheight}{14pt}
\setlength{\parindent}{0pt}
\setlength{\parskip}{6pt}
\setlength{\emergencystretch}{2em}
\setlist{itemsep=3pt,topsep=4pt}
\lstset{language=Python,basicstyle=\small\ttfamily,breaklines=true,
 columns=fullflexible,keepspaces=true,showstringspaces=false,
 frame=single,rulecolor=\color{gray!35},backgroundcolor=\color{gray!4},
 commentstyle=\color{green!35!black},keywordstyle=\color{blue!65!black},
 aboveskip=8pt,belowskip=8pt}
\newcommand{\code}[1]{\texttt{\detokenize{#1}}}
\newcommand{\RSS}{\operatorname{RSS}}
\newcommand{\RMSE}{\operatorname{RMSE}}
\newcommand{\MAE}{\operatorname{MAE}}
\newcommand{\E}{\mathbb{E}}
\newcommand{\ind}[1]{\mathbf{1}\{#1\}}
\newcommand{\fig}[2]{\begin{figure}[htbp]\centering
\includegraphics[width=\linewidth,height=0.39\textheight,keepaspectratio]{figures/#1}
\caption{#2}\end{figure}}
\newcommand{\notebooktitle}[3]{\begin{center}{\LARGE\bfseries #1\par}\vspace{5pt}
{\large Problem statement and worked resolution\par}\vspace{5pt}
Companion to \code{#2}\par\end{center}
\textbf{How to use this explanation.} #3 Section numbers follow the notebook's numbered sections where present. Numerical values are rounded and refer to the recorded run, not universal performance claims.}
\newcommand{\methods}[1]{\subsection*{Main Python calls}#1}

\begin{document}
\notebooktitle{3B: Kernels, curved boundaries, and scaling}{03_B_kernels_and_scaling.ipynb}{This resolution explains the feature-space equivalence, why linear models struggle with the chosen target, and how to interpret the kernel search without selecting from test results.}
\section*{Problem statement}
Generate 1,200 independent uniform points on $[-2,2]^2$. The clean label is positive inside the circle
$x_1^2+x_2^2<1.5$ and negative outside. Independently flip each label with probability 0.08.
Fit and evaluate these noisy observed labels. The clean boundary is available only because this is a simulation.
A learner is being asked to predict observed outcomes, not an easier clean target substituted during scoring.

A kernel $K(x,z)$ is an inner product between transformed features. A support vector machine fits a linear
rule in that feature space, which may have a nonlinear boundary in the original coordinates.
C balances hinge loss against coefficient norm; RBF gamma controls the distance scale of similarity.
Both must be chosen using development data rather than the final test.

\section{Generate data and reserve independent roles}
Simulation seed 17 and stratified split seed 23 give 720 training, 240 validation, and 240 test rows.
Training has about 31\% positive labels, so accuracy alone rewards a negative-majority predictor.
Balanced accuracy, the mean of the two class recalls, gives that constant predictor a score of 0.5.
A pipeline learns scaling only from fitting rows and reuses that transformation for evaluation.

The circle is centered on the origin. For the ideal population distribution, neither $x_1$ nor $x_2$ has
a simple signed linear association with being inside: points with large positive and large negative coordinates
are both likely outside. Squared coordinates carry the relevant information. This is a representation issue,
not evidence that the variables themselves are irrelevant.

\section{Derive and check the quadratic feature map}
Expand the kernel:
\[
(1+x^Tz)^2=1+2x_1z_1+2x_2z_2+x_1^2z_1^2+2x_1x_2z_1z_2+x_2^2z_2^2.
\]
It equals the inner product of
\[\phi(x)=(1,\sqrt2x_1,\sqrt2x_2,x_1^2,\sqrt2x_1x_2,x_2^2).\]
The square-root factors make each squared coefficient in the dot product equal the required factor 2.
A map with unweighted columns spans the same polynomial functions but gives them a different Euclidean norm,
so the same numeric C would impose a different regularization preference.

The notebook first standardizes the original two inputs using training statistics, then applies either the map
or the matching kernel. It does not standardize mapped coordinates separately. This distinction is essential:
rescaling the expanded columns would change the kernel identity being tested.
\code{FunctionTransformer} later wraps the map for use in a pipeline; it computes fixed columns rather than
learning a new statistical transformation.

With C=1 and tight tolerance, the validation-score difference is about $2.6\times10^{-13}$.
Both primal objectives are about 237.499742. For the implicit fit, the squared feature-space norm is recovered as
\[\|w_\phi\|^2=c^TK_{SV}c,\qquad c_i=\alpha_i y_i.\]
This makes objective comparison possible without explicitly recovering an infinite-dimensional or hidden weight vector.
The constant map coordinate does not require a second effective intercept: any constant contribution can be
moved into the unpenalized intercept, reducing its penalized coefficient to zero at an optimum.

\subsection*{Why prediction agreement is a stronger check than a matching label table}
Two models can make the same signs while assigning very different scores and optimizing different penalties.
Here the check compares the Gram matrix, real-valued scores, and objective values under matched conventions.
Numerical tolerance is appropriate because solvers and floating-point kernel calculations can differ slightly.

\section{Choose settings within and across model families}
Search C in $\{0.01,0.1,1,10\}$. For RBF also search gamma in $\{0.01,0.1,1,10\}$.
This focused grid avoids very slow large-C linear fits near degeneracy, and uses the ordinary SVC tolerance of 0.001
for model selection. The separate equivalence check retains its tighter tolerance.
The grid is a teaching choice, not a claim that stronger penalties could never be useful.

Each of five training folds fits preprocessing and its model on 576 rows and evaluates on 144.
Logistic regression uses its own loss and penalty convention, even though its inverse-regularization parameter
also happens to be named C. The explicit and implicit quadratic SVMs use the same feature-space objective.
\begin{center}\begin{tabular}{lrrr}\toprule
Family&Selected C&Gamma where relevant&CV balanced accuracy\\\midrule
Logistic&0.01&---&0.500000\\
Linear SVM&0.01&---&0.500000\\
Explicit quadratic&0.1&---&0.899253\\
Polynomial&0.1&1&0.899253\\
RBF&10&0.1&0.897455\\\bottomrule\end{tabular}\end{center}
The explicit quadratic and polynomial models tie; table order selects the explicit representation as the overall family.
This does not mean the equivalent polynomial rule is inferior. The slightly lower RBF CV value is selection evidence,
not a precise ranking of the population errors. The chosen setting in each family is refitted on all training rows.
\fig{03_B_kernels_and_scaling_1.png}{The RBF grid displays held-out balanced accuracy, not training accuracy. Both C and gamma influence fit flexibility.}

\section{Interpret boundaries and the change-of-units diagnostic}
The model with tiny C and gamma has broad similarities and strong pressure toward a low-norm fit.
At the opposite fixed diagnostic setting, C=100 and gamma=10 can form much more local boundaries.
The support-vector circles identify training rows entering the decision rule, not only mislabeled or difficult rows.
The clean circle is drawn for explanation, not for selecting a model from test performance.
\fig{03_B_kernels_and_scaling_2.png}{Boundary panels for selected models and two prespecified RBF diagnostics. A smooth curved rule can capture the circle without matching every randomly flipped label.}

Multiply only the first coordinate by 100. Raw squared distance becomes
$10000(x_1-z_1)^2+(x_2-z_2)^2$. Adjusting one scalar gamma cannot undo this uneven change for every pair.
With fixed selected RBF settings, raw validation balanced accuracy changes from about 0.8746 to 0.7296.
The standardized original and rescaled fits both score about 0.8715, and their decision scores agree within tolerance.

Why does standardization restore agreement? For a positive unit multiplier $a$, the transformed mean and scale
are $a\mu$ and $a\sigma$, so $(ax-a\mu)/(a\sigma)=(x-\mu)/\sigma$.
The equivalence assumes the scaler is refitted on the correspondingly rescaled training inputs and then reused.
Fitting a fresh scaler only on validation would be a different, invalid operation.
The raw original-unit fit happens to score slightly higher here; the purpose of scaling is to control geometry,
not to guarantee that every finite validation score improves.

\section{Evaluate the frozen models and explain the ranking}
\begin{center}\begin{tabular}{lrrrr}\toprule
Model&Accuracy&Balanced accuracy&ROC AUC&Support vectors\\\midrule
Majority&0.6917&0.5000&0.5000&---\\
Logistic&0.6917&0.5000&0.5033&---\\
Linear SVM&0.6917&0.5000&0.4639&447\\
Explicit quadratic&0.8958&0.8798&0.9177&285\\
Polynomial&0.8958&0.8798&0.9177&285\\
RBF&0.9083&0.8925&0.9171&247\\\bottomrule\end{tabular}\end{center}
The RBF model has higher test accuracy than the CV-selected quadratic model, but the selection remains frozen.
Its AUC is slightly lower: the metric assesses score ordering rather than decisions at a particular cutoff.
Linear models largely predict the majority class at their zero cutoffs and fail to describe the radial structure.
The many support vectors in the linear SVM do not rescue its inappropriate representation.

The clean rule has 92\% expected ordinary accuracy against the independently flipped labels. This is not a
fixed upper bound on a finite test accuracy, nor the same number as optimal balanced accuracy, whose class
weighting differs. A test result near 90\% is compatible with both finite-sample model error and unavoidable
flips; reporting clean-label accuracy instead would change the question.

\newpage
\section*{Documentation and practical tips}
The \href{https://scikit-learn.org/1.8/modules/svm.html}{official SVM guide} explains kernel options and
support-vector prediction. In a polynomial SVC, \code{coef0} belongs to the kernel expression; \code{intercept_}
is the fitted decision offset. With a pipeline, support vectors are stored in transformed coordinates,
while \code{support_} can index the original fitting rows for a plot.

\begin{itemize}
\item Keep the exact feature map and kernel constants together when verifying equivalence.
\item A constant baseline's ROC AUC is 0.5 even if class imbalance gives much higher accuracy.
\item Select gamma jointly with C and fix preprocessing first; neither a training plot nor support count alone is a selection metric.
\item Extreme C can slow optimization without improving held-out performance. A focused grid and documented tolerances are preferable to silently accepting a nonconverged fit.
\item Preserve the pipeline: saving only support vectors without the input transformation does not preserve the original prediction function.
\end{itemize}
\end{document}
